\documentclass[11pt,a4paper]{article}

% ---------- Idioma y codificación ----------
\usepackage[spanish,es-tabla]{babel}
\usepackage[utf8]{inputenc}
\usepackage[T1]{fontenc}
\usepackage{lmodern}

% ---------- Formato de página ----------
\usepackage[margin=2.5cm]{geometry}
\usepackage{setspace}
\setstretch{1.15}
\usepackage{parskip}

% ---------- Tablas y listas ----------
\usepackage{booktabs}
\usepackage{tabularx}
\usepackage{array}
\usepackage{enumitem}

% ---------- Títulos en negrita (sin color) ----------
\usepackage{titlesec}
\titleformat{\section}
  {\Large\bfseries}
  {\thesection}{0.8em}{}
\titleformat{\subsection}
  {\large\bfseries}
  {\thesubsection}{0.6em}{}
\titleformat{\subsubsection}
  {\normalsize\bfseries}
  {\thesubsubsection}{0.5em}{}
\titlespacing*{\section}{0pt}{18pt}{8pt}

% ---------- Encabezado y pie ----------
\usepackage{fancyhdr}
\pagestyle{fancy}
\fancyhf{}
\fancyhead[L]{\small\textbf{Scrum --- Sprint \arabic{sprintnum}}}
\fancyhead[R]{\small\textbf{CrediStep}}
\fancyfoot[C]{\small\thepage}
\renewcommand{\headrulewidth}{0.4pt}

% ---------- Número de sprint ----------
\newcounter{sprintnum}
\setcounter{sprintnum}{2}   % Cambie este número en cada sprint

% ---------- Caja del objetivo (solo blanco y negro) ----------
\usepackage[most]{tcolorbox}
\newtcolorbox{objetivo}{
  colback=white, colframe=black,
  boxrule=0.8pt, arc=0pt,
  left=8pt, right=8pt, top=6pt, bottom=6pt,
  fonttitle=\bfseries, title={Objetivo del Sprint}
}

% ======================================================
\begin{document}

% ---------- Portada ----------
\begin{titlepage}
    \centering
    \begin{figure}
        \centering
        \includegraphics[width=0.5\linewidth]{1.png}
        \label{fig:placeholder}
    \end{figure}
    {\Large \textbf{Pontificia Universidad Javeriana}} \\
    \vspace{0.2cm}
    {\large Ingeniería de Sistemas} \\
    {\large Gestión de Innovación en TI} \\
    
    \vspace{3.5cm}
    
    {\Huge \textbf{CrediStep}} \\
    \vspace{0.4cm}
    {\Large Documento Scrum — Sprint 2} \\
    \vspace{0.2cm}
    {\large Elaboración y presentación del Pitch } \\
    
    \vfill
    
    \textbf{Integrantes:} \\
    \vspace{0.2cm}
    Beiker Romario Jaramillo Giraldo \\
    Laura Sofía Ortiz Gómez \\
    Sarahy Michelle Silva Pinto \\
    Valentina Carrillo Peñuela \\
    Miguel Ángel Gómez López \\
    
    \vspace{1cm}
    
    \textbf{Profesor:} Víctor Hugo Vela Pinzón \\
    \vspace{0.8cm}
    Bogotá D.C. \\
    Inicio: 03/08/2026 \\
    Fin: 31/08/2026
    
\end{titlepage}

\section{Información general}

\begin{tabularx}{\textwidth}{>{\bfseries}l X}
  \toprule
  Sprint               & \arabic{sprintnum} \\
  Proyecto             & CrediStep \\
  Product Owner        & Víctor Hugo Vela Pinzón (profesor) \\
  Scrum Master         & Beiker Romario Jaramillo Giraldo \\
  Equipo de desarrollo & Laura Sofía Ortiz Gómez, Sarahy Michelle Silva Pinto, Valentina Carrillo Peñuela, Miguel Ángel Gómez López \\
  Fecha de inicio      & 03/08/2026 \\
  Fecha de fin         & 31/08/2026 \\
  Duración             & 4 semanas \\
  Capacidad del equipo & 10 puntos \\
  \bottomrule
\end{tabularx}

% ---------- 2. Objetivo del sprint ----------
\section{Objetivo del sprint}

\begin{objetivo}
  Elaborar y presentar un pitch de nueve partes sobre CrediStep, con diapositivas y guion que cuenten la misma historia.
\end{objetivo}

% ---------- 3. Backlog del sprint ----------
\section{Backlog del sprint}

\subsection{Historias de usuario seleccionadas}

\begin{tabularx}{\textwidth}{c X c c c}
  \toprule
  \textbf{ID} & \textbf{Historia de usuario} & \textbf{Prioridad} & \textbf{Puntos} & \textbf{Estado} \\
  \midrule
  HU-01 & Como jurado del pitch, quiero nueve diapositivas que vayan del problema a la solución, para seguir el pitch con facilidad. & Alta & 5 & Terminada \\
  HU-02 & Como joven oyente del pitch, quiero un guion con las nueve partes y un tono cercano, para entender la idea sin tecnicismos. & Alta & 3 & Terminada \\
  HU-03 & Como Product Owner, quiero que cada dato del pitch traiga su fuente y año, para poder verificar las cifras. & Media & 2 & Terminada \\
  \bottomrule
\end{tabularx}

% ---------- 4. Desglose de tareas ----------
\section{Desglose de tareas}

\begin{tabularx}{\textwidth}{c X l c c}
  \toprule
  \textbf{HU} & \textbf{Tarea} & \textbf{Responsable} & \textbf{Horas est.} & \textbf{Estado} \\
  \midrule
  HU-01 & Diseñar las nueve diapositivas & Valentina Carrillo & 5 & Terminada \\
  HU-02 & Escribir el guion de las nueve partes & Sarahy Silva & 4 & Terminada \\
  HU-02 & Ensayar el pitch y alinear guion con diapositivas & Beiker Jaramillo & 3 & Terminada \\
  HU-03 & Reunir los datos y las fuentes del pitch & Laura Ortiz & 2 & Terminada \\
  HU-03 & Revisar las cifras de las diapositivas & Miguel Gómez & 2 & Terminada \\
  \bottomrule
\end{tabularx}

% ---------- 5. Seguimiento diario (Daily Scrum) ----------
\section{Seguimiento diario (Daily Scrum)}

\begin{tabularx}{\textwidth}{c l X X X}
  \toprule
  \textbf{Día} & \textbf{Miembro} & \textbf{Ayer} & \textbf{Hoy} & \textbf{Impedimentos} \\
  \midrule
  1 & Sarahy Silva & Definimos la estructura del pitch & Escribir la primera versión del guion & Ninguno \\
  2 & Beiker Jaramillo & Revisamos el guion con las diapositivas & Ensayar el pitch completo & Los títulos del guion y de las diapositivas no coinciden \\
  \bottomrule
\end{tabularx}

\subsection{Impedimentos y acciones}

\begin{tabularx}{\textwidth}{X l l}
  \toprule
  \textbf{Impedimento} & \textbf{Responsable} & \textbf{Estado} \\
  \midrule
  Los títulos del guion y de las diapositivas no coincidían & Beiker Jaramillo & Resuelto \\
  El deck usado en el pitch presencial no se envió al profesor & Beiker Jaramillo & Abierto \\
  \bottomrule
\end{tabularx}

% ---------- 6. Métricas del sprint ----------
\section{Métricas del sprint}

\begin{tabularx}{\textwidth}{X c c}
  \toprule
  \textbf{Indicador} & \textbf{Planeado} & \textbf{Real} \\
  \midrule
  Puntos de historia completados & 10 & 10 \\
  Historias completadas          & 3 & 3 \\
  Velocidad del equipo           & 10 & 10 \\
  \bottomrule
\end{tabularx}

% ---------- 7. Revisión del sprint ----------
\section{Revisión del sprint}

\subsection{Incrementos entregados}
\begin{itemize}[leftmargin=1.5em]
  \item {Pitch de nueve diapositivas presentado en clase.}
  \item {Guion del pitch con nueve partes.}
\end{itemize}

\subsection{Retroalimentación de los interesados}
{Comentarios del profesor (nota: 38/50, equivalente a 3,8):}
\begin{itemize}[leftmargin=1.5em]
  \item {La actividad breve de inicio estuvo bien; hay que saber manejarla y practicarla para que aporte valor al pitch.}
  \item {Los indicadores de contexto y oportunidad estuvieron bien.}
  \item {El perfilamiento del usuario estuvo bien; pidió profundizarlo en la validación.}
  \item {Nos pasamos del tiempo: usamos casi el doble de lo esperado.}
  \item {La propuesta de solución en cuatro pasos estuvo bien.}
  \item {En «¿Por qué nosotros?» faltó presentar al equipo y un buen cierre.}
  \item {No enviamos el deck (diapositivas) usado en el pitch presencial.}
\end{itemize}

\subsection{Historias no completadas}
\begin{tabularx}{\textwidth}{c X X}
  \toprule
  \textbf{ID} & \textbf{Motivo} & \textbf{Acción propuesta} \\
  \midrule
  -- & Todas las historias se completaron. & No aplica \\
  \bottomrule
\end{tabularx}

% ---------- 8. Retrospectiva ----------
\section{Retrospectiva del sprint}

\subsection{Qué salió bien}
\begin{itemize}[leftmargin=1.5em]
  \item {El pitch cubre las secciones esperadas y tiene un mensaje central fácil de recordar.}
  \item {El guion y las diapositivas cuentan la misma historia, con un tono cercano a los jóvenes.}
\end{itemize}

\subsection{Qué se puede mejorar}
\begin{itemize}[leftmargin=1.5em]
  \item {Usamos casi el doble del tiempo previsto.}
  \item {Faltó presentar al equipo y cerrar mejor.}
  \item {Las diapositivas con datos no muestran la fuente, y el deck no se envió.}
\end{itemize}

\subsection{Acciones de mejora}
\begin{tabularx}{\textwidth}{X l l}
  \toprule
  \textbf{Acción} & \textbf{Responsable} & \textbf{Fecha límite} \\
  \midrule
  Ensayar con cronómetro para cumplir el tiempo & Beiker Jaramillo & 21/09/2026 \\
  Agregar una diapositiva del equipo y reforzar el cierre & Valentina Carrillo & 21/09/2026 \\
  Citar la fuente y el año en cada diapositiva con datos & Miguel Gómez & 21/09/2026 \\
  Enviar el deck junto con cada pitch & Beiker Jaramillo & 21/09/2026 \\
  \bottomrule
\end{tabularx}


\end{document}
